\documentclass[10pt,a4paper]{article}

% =========================================================
% JOURNAL OF AGRONOMY AND CROP SCIENCE
% Manuscript template inspired by Wiley layout
%
% Required files in Overleaf:
% Journal.png
% original.png
% wiley.png
% ref.bib
% =========================================================

% =========================================================
% ENCODING / LANGUAGE
% =========================================================
\usepackage[T1]{fontenc}
\usepackage[utf8]{inputenc}
\usepackage[english]{babel}

% =========================================================
% FONT / TYPOGRAPHY
% =========================================================
\usepackage{newtxtext,newtxmath}
\usepackage{microtype}
\usepackage{xurl}
\usepackage{graphicx}%
\usepackage{multirow}%
\let\Bbbk\relax
\usepackage{amsmath,amssymb,amsfonts}%
\let\openbox\relax
\usepackage{amsthm}%
\usepackage{mathrsfs}%
\usepackage[title]{appendix}%
\usepackage{xcolor}%
\usepackage{textcomp}%
\usepackage{manyfoot}%
\usepackage{booktabs}%
\usepackage{algorithm}%
\usepackage{algorithmicx}%
\usepackage{algpseudocode}%
\usepackage{listings}%

% For orcid links generation
\RequirePackage{orcidlink}


% =========================================================
% GRAPHICS / COLORS
% =========================================================
\usepackage{graphicx}
\usepackage{xcolor}
% =========================================================
% PAGE GEOMETRY
% =========================================================
\usepackage[
    a4paper,
    top=1.35cm,
    bottom=1.45cm,
    left=1.60cm,
    right=1.60cm,
    headheight=14pt,
    headsep=0.25cm,
    footskip=0.60cm
]{geometry}

% =========================================================
% TEXT / LAYOUT
% =========================================================
\usepackage{ragged2e}
\usepackage{setspace}
\usepackage{titlesec}
\usepackage{enumitem}
\usepackage{calc}

% =========================================================
% TABLES / FIGURES
% =========================================================
\usepackage{booktabs}
\usepackage{multirow}
\usepackage{tabularx}
\usepackage{array}
\usepackage{caption}
\usepackage{dblfloatfix}

% =========================================================
% MATH
% =========================================================
\usepackage{amsmath}

\usepackage{amsfonts}

% =========================================================
% HEADER / FOOTER
% =========================================================
\usepackage{fancyhdr}
\usepackage{lastpage}

% =========================================================
% LINKS
% =========================================================
\usepackage{hyperref}

% =========================================================
% BIBLIOGRAPHY
% =========================================================
\usepackage[
  authordate,
  backend=biber,
  natbib=true,
  doi=true,
  url=false,
  isbn=false,
  giveninits=true,
  uniquename=false,
  maxcitenames=2,
  maxbibnames=99
]{biblatex-chicago}

\defbibenvironment{bibliography}
  {\list{}
    {\setlength{\leftmargin}{0pt}%
     \setlength{\itemindent}{0pt}%
     \setlength{\itemsep}{\bibitemsep}%
     \setlength{\parsep}{\bibparsep}}}
  {\endlist}
  {\item}

\DeclareFieldFormat{doi}{%
  \href{https://doi.org/#1}{https://doi.org/#1}%
}

\addbibresource{ref.bib}

% =========================================================
% GRAPHIC PATHS
% =========================================================
\graphicspath{{./}{figures/}{images/}}

% =========================================================
% GENERAL CONFIGURATION
% =========================================================
\setlength{\parindent}{0pt}
\setlength{\parskip}{0.28em}

\setlength{\columnsep}{0.65cm}
\setlength{\columnseprule}{0pt}

\setlength{\emergencystretch}{2em}

\linespread{0.98}

\raggedbottom

% =========================================================
% COLORS
% =========================================================
\definecolor{JACGray}{RGB}{230,230,230}
\definecolor{JACDark}{RGB}{36,32,33}
\definecolor{JACBlue}{RGB}{0,92,165}
\definecolor{JACRule}{RGB}{75,75,75}

% =========================================================
% HYPERREF
% =========================================================
\hypersetup{
    colorlinks=true,
    linkcolor=JACBlue,
    citecolor=JACBlue,
    urlcolor=JACBlue,
    pdfauthor={Fabrício Correia de Oliveira},
    pdftitle={Hybrid crop modelling and machine learning for second-season maize in Paraná}
}

% =========================================================
% SECTION TITLES
% =========================================================
\titleformat{\section}
    {\large\bfseries\color{JACDark}}
    {\thesection\,|}
    {0.40em}
    {}

\titleformat{\subsection}
    {\normalsize\bfseries\color{JACDark}}
    {\thesubsection\,|}
    {0.40em}
    {}

\titleformat{\subsubsection}
    {\normalsize\bfseries\color{JACDark}}
    {\thesubsubsection\,|}
    {0.40em}
    {}

\titlespacing*{\section}
    {0pt}
    {0.85em}
    {0.25em}

\titlespacing*{\subsection}
    {0pt}
    {0.60em}
    {0.15em}

\titlespacing*{\subsubsection}
    {0pt}
    {0.50em}
    {0.10em}

% =========================================================
% CAPTIONS
% =========================================================
\captionsetup{
    font=small,
    labelfont=bf,
    labelsep=quad,
    justification=justified,
    singlelinecheck=false
}

% =========================================================
% LISTS
% =========================================================
\setlist[itemize]{
    leftmargin=1.2em,
    itemsep=0.12em,
    topsep=0.15em,
    parsep=0pt
}

\setlist[enumerate]{
    leftmargin=1.4em,
    itemsep=0.12em,
    topsep=0.15em,
    parsep=0pt
}

% =========================================================
% HEADER / FOOTER — PAGE 2 ONWARD
% =========================================================
\pagestyle{fancy}

\fancyhf{}

\renewcommand{\headrulewidth}{0pt}
\renewcommand{\footrulewidth}{0pt}

\fancyhead[L]{%
    \small\itshape
    Journal of Agronomy and Crop Science
}

\fancyhead[R]{%
    \small
    \thepage\ of \pageref{LastPage}
}

\fancyfoot[L]{%
    \small\itshape
    Journal of Agronomy and Crop Science
}

\fancyfoot[R]{%
    \small
    \thepage\ of \pageref{LastPage}
}

% =========================================================
% ARTICLE TITLE
% =========================================================
% USER INPUT: replace the placeholder text below.
% Journal guidance: use a brief, informative title containing the major keywords.
% Avoid abbreviations in the title.
\newcommand{\ArticleTitle}{%
Hybrid crop modelling and machine learning for second-season maize yield and the economic risk of sowing dates in Paraná, Brazil
}

% =========================================================
% AUTHORS
% =========================================================
% Add or remove authors as needed.
% Enter the full names of all authors.
% Use superscript numbers to link authors to affiliations.
% Use * for the corresponding author.
% The journal requires ORCID identifiers.
% Example ORCID format: \orcid{0000-0000-0000-0000}

\newcommand{\orcid}[1]{%
  \hspace{1pt}\orcidlink{#1}\hspace{1pt}%
}

\newcommand{\ArticleAuthors}{%
Fabrício Correia de Oliveira\orcid{0000-0000-0000-0000}\textsuperscript{1,*}
\;|\;
[INSERT: co-authors and ORCID identifiers]
}


% =========================================================
% AFFILIATIONS
% =========================================================
% Add, remove, or renumber affiliations as needed.
% Include the institutional affiliation where the work was conducted.
% If an author's present address differs, provide it as a footnote or note as appropriate.
\newcommand{\AffiliationsBlock}{%

\textsuperscript{1}
Universidade Tecnológica Federal do Paraná (UTFPR), Campus Santa Helena, Santa Helena, Paraná, [INSERT: postal code], Brazil.

}


% =========================================================
% CORRESPONDENCE
% =========================================================
\newcommand{\CorrespondenceBlock}{%
Correspondence:
Fabrício Correia de Oliveira
(\href{mailto:fcoliveira@utfpr.edu.br}{fcoliveira@utfpr.edu.br})
}


% =========================================================
% EDITORIAL DATES
% =========================================================
\newcommand{\DatesBlock}{%
Received: -- \;|\;
Revised: -- \;|\;
Accepted: --
}

% =========================================================
% KEYWORDS
% =========================================================
% Journal guidance: provide up to five keywords.
% Include at least one relevant stress topic when applicable to the manuscript.
% Avoid unnecessary duplication of words already used in the title.
\newcommand{\KeywordsBlock}{%
agro-ecological zone model \;|\;
crop insurance \;|\;
drought and heat stress \;|\;
risk analysis \;|\;
\textit{Zea mays}
}


% =========================================================
% ABBREVIATIONS
% =========================================================
% If no abbreviations are used, remove this content from the first-page footer.
\newcommand{\AbbreviationsBlock}{%
APSIM, Agricultural Production Systems sIMulator; CE, certainty equivalent; CVaR, conditional value at
risk; DSSAT, Decision Support System for Agrotechnology Transfer; DTA, total available soil water;
ISNA, water requirement satisfaction index; MZA-FAO, FAO Agro-Ecological Zone model; PA, attainable
yield; PP$_f$, potential yield; PSR, federal crop insurance subsidy programme; SERF, stochastic efficiency
with respect to a function; SHAP, Shapley additive explanations; ZARC, Agricultural Climate Risk Zoning.
}


% =========================================================
% JOURNAL INFORMATION
% =========================================================
\newcommand{\RunningYear}{20XX}

\newcommand{\VolumeInfo}{%
Volume:ArticleNumber
}

\newcommand{\DOIInfo}{%
https://doi.org/xx.xxxx/xxxxx
}

% =========================================================
% ABSTRACT
% =========================================================
% Journal guidance: the abstract should be unstructured and should not contain abbreviations.
% Briefly state the context or problem, objective, main methods, principal results, and conclusion.
% Keep the text concise and focused on the study's main contribution.
\newcommand{\AbstractText}{%
\noindent
Second-season maize, sown after soybean, is the main source of maize in Brazil, and its sowing date
determines how much the crop is exposed to drought, frost and heat. This study developed and validated a
hybrid model of second-season maize yield for the 399 municipalities of Paraná and established the best
sowing date of each municipality from its economic risk. Four process-based crop models of increasing
complexity, used without calibration, simulated 27 seasons and nine sowing dates with daily gridded
weather and generated the features of machine learning models trained on official municipal yields and
validated by leaving whole years out. A simple agro-ecological zone model was as informative as the more
detailed models, and its ensemble of a physically constrained gradient boosting and a support vector
regression explained 43\% of the variance of municipal yields and 67\% of the variation of the state mean
between years. Water satisfaction at flowering and grain filling, reproductive frost and heat at
flowering were the main drivers of the yield anomaly. Combining the yield distribution with the prices,
costs and insurance access of each season, the certainty equivalent of the margin was 2.5 times higher for
sowing on 5 January than on 25 February, and early sowing was the best option in 84\% of the
municipalities; only the warm lowlands of the west and northwest preferred mid-March, to avoid heat at
flowering. The modelled frequency of crop failures followed the indemnified insurance policies between
years, and the recommendation was robust to costs, prices, insurance terms and crop model parameters.
Sowing as early as the soybean harvest allows is the most profitable and least risky strategy for
second-season maize in most of Paraná.
}


% =========================================================
% ABSTRACT BOX
% =========================================================
\newcommand{\AbstractBox}[1]{%

\vspace{0.30cm}

\noindent
\colorbox{JACGray}{%
    \parbox{\dimexpr\textwidth-2\fboxsep\relax}{%

        \vspace{0.20cm}

        {\large\bfseries ABSTRACT}

        \vspace{0.10cm}

        \justifying

        #1

        \vspace{0.14cm}

    }%
}

}

% =========================================================
% DOCUMENT
% =========================================================
\begin{document}

% =========================================================
% FIRST PAGE
% =========================================================
\thispagestyle{empty}

% =========================================================
% JOURNAL HEADER
% =========================================================
\noindent
\begin{minipage}[t]{0.45\textwidth}

\vspace{0pt}

\includegraphics[
    width=0.90\linewidth,
    keepaspectratio
]{Journal.png}

\end{minipage}
%
\hfill
%
\begin{minipage}[t]{0.27\textwidth}

\vspace{-0.10cm}

\raggedleft

\includegraphics[
    width=0.95\linewidth,
    keepaspectratio
]{wiley.png}

\end{minipage}

\vspace{0.52cm}

% =========================================================
% ORIGINAL ARTICLE / OPEN ACCESS
% =========================================================
\noindent

\includegraphics[
    width=0.40\textwidth,
    keepaspectratio
]{original.png}

\vspace{0.48cm}

% =========================================================
% TITLE
% =========================================================
{\fontsize{21}{24}\selectfont
\bfseries
\color{JACDark}
\ArticleTitle
\par}

\vspace{0.28cm}

% =========================================================
% AUTHORS
% =========================================================
{\small
\ArticleAuthors
\par}

\vspace{0.18cm}

% =========================================================
% AFFILIATIONS
% =========================================================
{\footnotesize

\setlength{\parindent}{0pt}

\raggedright

\AffiliationsBlock

\par
}

\vspace{0.18cm}

% =========================================================
% CORRESPONDENCE
% =========================================================
{\footnotesize
\CorrespondenceBlock
\par
}

\vspace{0.08cm}

% =========================================================
% DATES
% =========================================================
{\footnotesize

\textbf{Received:} -- 
\;|\;
\textbf{Revised:} --
\;|\;
\textbf{Accepted:} --

\par
}

\vspace{0.08cm}

% =========================================================
% KEYWORDS
% =========================================================
{\footnotesize

\textbf{Keywords:}
\KeywordsBlock

\par
}

% =========================================================
% ABSTRACT
% =========================================================
\AbstractBox{\AbstractText}

\vspace*{\fill}

% =========================================================
% FIRST-PAGE FOOTER
% =========================================================
{\scriptsize

\hrule

\vspace{0.10cm}

\textbf{Abbreviations:}
\AbbreviationsBlock

\vspace{0.10cm}

\hrule

\vspace{0.10cm}

This is an open access article under the terms of the
Creative Commons Attribution License, which permits use,
distribution and reproduction in any medium,
provided the original work is properly cited.

\vspace{0.08cm}

© \RunningYear\ The Author(s).
\textit{Journal of Agronomy and Crop Science}
published by Wiley-VCH GmbH.

\vspace{0.10cm}

\hrule

\vspace{0.10cm}

\textit{Journal of Agronomy and Crop Science},
\RunningYear;
\VolumeInfo

\href{\DOIInfo}{\DOIInfo}

\hfill
1 of \pageref{LastPage}

}

% =========================================================
% PAGE 2 ONWARD — TWO COLUMNS
% =========================================================
\clearpage

\twocolumn

\pagestyle{fancy}

% =========================================================
% KEY POINTS
% =========================================================
% Journal guidance: provide 3 to 5 brief Key Points in bullet form.
% Each point should summarize the initial situation/problem, approach, main result,
% significance, or implication of the research.
% Keep each point brief (approximately 50--100 characters) and avoid abbreviations.
% Replace the example items below with the main take-home messages of the manuscript.
\noindent
\fbox{%
\begin{minipage}{0.94\columnwidth}

\vspace{0.14cm}

{\normalsize\bfseries Key Points}

\vspace{0.06cm}

\begin{itemize}
    \setlength{\itemsep}{0.05cm}
    \setlength{\topsep}{0.02cm}
    \setlength{\parsep}{0pt}
    \setlength{\partopsep}{0pt}

    \item Insert the first key point here.
    \item Insert the second key point here.
    \item Insert the third key point here.
    \item Insert the fourth key point here.
    \item Insert the fifth key point here.

\end{itemize}

\vspace{0.04cm}

\end{minipage}%
}

\vspace{0.22cm}

% =========================================================
% 1 INTRODUCTION
% =========================================================
\section{Introduction}
\label{sec:introduction}

Among agricultural activities, crop production is the most dependent on weather conditions: climate
variation explains about one third of the global variability of crop yields. In Brazil this dependence
is particularly evident in second-season maize, sown immediately after the soybean harvest in a
double-crop system. This crop has become the main source of Brazilian maize,
with about 101 of the 128 million tonnes produced in the 2024/25 season, and in Paraná, the second
largest producing state, it reached 17.3 million tonnes on 2.8 million hectares in 2025, 85\% of the
state's maize production. The expansion of the second season has also reshaped the seasonal pattern
of domestic maize prices, since most of the harvest is now concentrated in the middle of the year
\citep{ray2015, conab2025, ibge2025, souza2021}.

Because second-season maize is sown at the end of summer, its cycle advances into autumn and winter,
when rainfall, solar radiation and temperature decline. Water deficit is the main cause of the yield
gap of this crop in Brazil, and in Paraná, where the cycle extends into the coldest months, frost
during flowering and grain filling adds to it, particularly in the south-central and western regions.
Each day of delay after the soybean harvest therefore shifts the sensitive phases to less favourable
conditions, which makes the sowing date the most important management decision of the season; poorly
chosen sowing dates remain among the main causes of maize yield losses in Brazil, and the choice of the
sowing date alone can reduce the water-limited yield gap of off-season maize by 66--89\%. In Brazil, the Agricultural Climate Risk
Zoning (ZARC) indicates, for each municipality, soil water class and cultivar cycle, the ten-day sowing
periods with up to 40\% risk of yield loss, and planting within these periods is a condition for
subsidized crop insurance and for the PROAGRO programme. The ZARC, however, classifies climatic risk
and does not quantify the economic return of each date nor the risk attitude of the farmer
\citep{goncalves2002, andrea2018, noia2019eja, noia2019agsy, noia2020, monteiro2021, embrapa2023}.

Process-based crop models are the classical tool to quantify how weather and sowing date affect yield,
and they differ widely in complexity. The FAO Agro-Ecological Zone model (MZA-FAO) estimates the
potential yield from radiation and temperature and reduces it to the attainable yield through the
relative water deficit of each phenological phase, weighted by the yield response factor ($k_y$) and
computed from a crop water balance; it requires few inputs and parameters, works with gridded weather
data, and has been applied in Brazil to cultivar drought tolerance, yield gap analysis and operational
yield monitoring. AquaCrop, also developed by FAO, is water-driven,
simulates canopy cover and transpiration and converts transpiration into biomass through a normalized
water productivity and into yield through a harvest index. CSM-CERES-Maize, in the Decision Support
System for Agrotechnology Transfer (DSSAT), simulates at a daily step phenology, leaf area expansion,
kernel number and the soil water and nitrogen balances, driven by cultivar-specific coefficients.
The Agricultural Production Systems sIMulator (APSIM) integrates crop, soil and management modules in
a modular framework, now in its Next Generation version. The more detailed models represent more
processes, but they depend on calibration with field data for each cultivar and site, which is rarely
available at the municipal scale, and different models diverge considerably in their responses to
the same weather, both for maize worldwide and for crops in southern Brazil \citep{doorenbos1979,
allen1998, battisti2015, duarte2019, noia2020, noia2024, steduto2009, kelly2021, jones2003,
holzworth2014, bassu2014, battisti2017, seidel2018}.

Machine learning has become widespread in crop yield prediction, including soybean in Brazil, because
it learns directly from observed yields and handles nonlinear and interacting effects. Statistical
models, however, have no
physical structure, so they extrapolate poorly beyond the conditions of the training data, and when
municipality-years of the same season are randomly split between training and validation, information
leaks between years and the accuracy is overestimated. Hybrid approaches address both limitations by
using the outputs of a process-based model as features of a machine learning model: this strategy
improved the evaluation of climate extremes on wheat yield in Australia with APSIM and the prediction
of maize yield in the US Corn Belt, and physical knowledge can also be imposed on the learning itself,
for instance as constraints on the direction of the response. Published hybrids, however, rely on
detailed and usually calibrated crop models. It remains unknown whether a complex crop model, used
without calibration, provides more information to the hybrid model than the simple MZA-FAO, and whether
the hybrid predictions keep a physically consistent response when applied to sowing dates other than
those that generated the observed yields \citep{vanklompenburg2020, schwalbert2020, vonbloh2023,
lobell2010, roberts2017, kapoor2023, feng2019, shahhosseini2021, karpatne2017}.

The same limitation applies to the economic side of the decision. Studies on sowing dates usually
compare the mean yield or the mean profitability of each date, including in the soybean--maize system
of Paraná, whereas the farmer decides under risk:
the revenue depends on the yield and on the price of the harvest month, and most costs are committed
at sowing and are not recovered in a crop failure. Stochastic efficiency with respect to a function
(SERF) ranks risky alternatives by their certainty equivalent over a range of risk aversion and is
suited to this problem. Combining the yield distribution of a hybrid model with the prices and costs of
the same seasons and with the insurance access defined by the ZARC makes it possible to translate
climatic risk into economic risk for each municipality and sowing date, but this integration has not
yet been carried out at the state scale nor checked against independent records of crop losses
\citep{noia2019eja, oligini2021, hardaker2004}.

In this regard, this study aimed to develop and validate a hybrid model of second-season maize yield
for the 399 municipalities of Paraná, to quantify whether process-based crop models of increasing
complexity, used without calibration, improve its predictions, and to establish the best sowing date
for each municipality from the economic risk of each date, considering the farmer's risk aversion and
the access to crop insurance under the ZARC.


% =========================================================
% 2 MATERIALS AND METHODS
% =========================================================
\section{Materials and Methods}
\label{sec:methods}

\subsection{Study area and workflow}
\label{subsec:area}

The study covered the 399 municipalities of the state of Paraná, southern Brazil (22.5--26.7\textdegree{}S,
48.0--54.6\textdegree{}W), where altitude ranges from sea level to more than 1,000 m and the climate varies
from humid subtropical with hot summers in the north and west to humid subtropical with mild summers on
the southern plateau. Second-season maize is sown from January to March after early soybean, and the
analysis was organized in four steps: (i) daily weather and soil data for each municipality; (ii)
simulation of each municipality, year and sowing date with four process-based crop models;
(iii) training and validation of hybrid yield models against municipal yields; and (iv) economic risk
analysis of each sowing date and choice of the best date per municipality. All steps were run in Python
3.13, and the code and data are available in a public repository (see Data Availability Statement).

\subsection{Weather, soil and yield data}
\label{subsec:data}

Daily weather from 1 December 1998 to 31 December 2025 was extracted at the seat of each municipality
from the Brazilian Daily Weather Gridded Data (BR-DWGD, version 3.2.4, 0.1\textdegree{} grid), through
Google Earth Engine: precipitation, maximum and minimum air temperature, relative humidity, global solar
radiation, wind speed at 2 m and FAO-56 Penman--Monteith reference evapotranspiration ($ET_o$). This
dataset interpolates surface observations of the national networks and performed well for maize yield
simulation in Brazil \citep{xavier2022, duarte2019, allen1998}.

The total available water of the soil (DTA, mm m$^{-1}$) of each municipality was the mean of the soil
available water capacity map of Brazil (1:250,000) produced by the National Water and Sanitation Agency
(ANA) and the Federal University of Paraná for the Irrigation Atlas, over the area of
temporary crops (soybean and other temporary crops,
classes 39 and 41 of MapBiomas Collection 10, year 2024), and not the value at the municipal seat;
municipalities without mapped temporary crops used the mean of their whole territory. For APSIM,
AquaCrop and DSSAT, clay and sand contents in ten 10-cm layers (0--100 cm) were obtained from MapBiomas
Soil Collection 3 over the same cropland mask \citep{ana_awc, souza2020}. For the official risk zoning, each municipality was assigned the
available-water class AD1--AD6 of the Brazilian zoning rules from its DTA (17, 44, 54, 263 and 21
municipalities in classes AD2 to AD6).

Municipal second-season maize yield (kg ha$^{-1}$) from 2003 to 2025 was obtained from the Municipal
Agricultural Production survey (PAM, table 839) of IBGE, indexed by the harvest year. Municipality-years
with less than 100 ha harvested (residual second season) and values identical to those of the previous
year (8.5\% of the records, which indicate repeated estimates rather than surveys) were removed, leaving
6,806 municipality-years in 373 municipalities \citep{ibge2025}.

\subsection{MZA-FAO crop model}
\label{subsec:mza}

The MZA-FAO was run with the \texttt{agrometeorologiapy} library (version 0.3.0) for each municipality,
sowing year from 1999 to 2025 and nine sowing dates (days 5, 15 and 25 of January, February and March),
totalling 96,957 simulated cycles. Phenology followed the accumulated growing degree-days ($GDD$, eq. 1),
with base temperature $T_b$ = 10 \textdegree{}C and a thermal time of 1,450 \textdegree{}C d from sowing to
physiological maturity; cycles that did not reach this sum within 300 days were flagged as incomplete
\citep{pereira2002}.
\begin{equation}
GDD =
\begin{cases}
T_{med} - T_b, & T_b < T_{min}\\[2pt]
\dfrac{(T_{max} - T_b)^2}{2\,(T_{max} - T_{min})}, & T_{min} \le T_b < T_{max}\\[6pt]
0, & T_{max} \le T_b
\end{cases}
\label{eq:gdd}
\end{equation}
where $T_{max}$, $T_{min}$ and $T_{med}$ are the maximum, minimum and mean daily air temperatures
(\textdegree{}C).

The fraction of the thermal time defined the crop coefficient ($K_c$) in the four FAO-56 stages
(fractions 0.17, 0.28, 0.33 and 0.22; $K_c$ = 0.3, 1.2 and 0.6 for the initial, mid-season and end
stages, with linear transitions), the five yield response periods of Doorenbos and Kassam
(establishment, vegetative, flowering, grain filling and maturity, fractions 0.15, 0.30, 0.12, 0.30 and
0.13) and the effective rooting depth $z$, which grew linearly from 0.10 m at sowing to 0.40 m at the end
of the vegetative period. The daily crop water balance followed Thornthwaite and Mather, with the crop
evapotranspiration ($ET_c$) as the demand and an available water capacity ($CAD$) that increased with
the roots (eqs. 2--6). The balance started 30 days before sowing, with the initial $K_c$ and root depth,
starting from a full profile, to establish the soil water storage at sowing \citep{allen1998,
doorenbos1979, thornthwaite1955}.
\begin{equation}
ET_{c,i} = K_{c,i}\, ET_{o,i}, \qquad CAD_i = z_i\, DTA
\label{eq:etc}
\end{equation}
\begin{equation}
ARM_i =
\begin{cases}
ARM_{i-1}\,\exp\!\left[(P_i - ET_{c,i})/CAD_i\right], & P_i - ET_{c,i} < 0\\[2pt]
\min\!\left(ARM_{i-1} + P_i - ET_{c,i},\; CAD_i\right), & P_i - ET_{c,i} \ge 0
\end{cases}
\label{eq:arm}
\end{equation}
\begin{equation}
ETR_i =
\begin{cases}
P_i + \lvert ARM_i - ARM_{i-1} \rvert, & P_i - ET_{c,i} < 0\\
ET_{c,i}, & P_i - ET_{c,i} \ge 0
\end{cases}
\label{eq:etr}
\end{equation}
\begin{equation}
DEF_i = ET_{c,i} - ETR_i, \qquad
EXC_i =
\begin{cases}
P_i - ET_{c,i} - (ARM_i - ARM_{i-1}), & ARM_i = CAD_i\\
0, & ARM_i < CAD_i
\end{cases}
\label{eq:def}
\end{equation}
\begin{equation}
ISNA_j = \frac{\sum_{i \in j} ETR_i}{\sum_{i \in j} ET_{c,i}}
\label{eq:isna}
\end{equation}
where $i$ is the day; $j$ is the yield response period; $ET_o$, $ET_c$, $ETR$, $DEF$ and $EXC$ are the
reference, crop and actual evapotranspiration, water deficit and water surplus (mm d$^{-1}$); $P$ is the
precipitation (mm d$^{-1}$); $ARM$ is the soil water storage (mm); $CAD$ is the available water capacity
(mm); $z$ is the effective rooting depth (m); $DTA$ is the total available water of the soil
(mm m$^{-1}$); and $ISNA$ is the water requirement satisfaction index (dimensionless).

The potential yield ($PP_f$) was computed day by day with the agro-ecological zone method, as the sum of
the gross biomass production of a standard crop under clear and overcast skies, corrected for leaf area,
respiration, harvest index and grain moisture (eqs. 7--9). Because BR-DWGD provides global radiation and
not sunshine duration, the relative sunshine $n/N$ was obtained by inverting the Angström--Prescott
equation and limited to the interval 0--1. The temperature corrections for C4 crops followed the polynomials of Barbieri and Tuon
\citep{doorenbos1979, barbieri1992, allen1998, battisti2015}.
\begin{equation}
\frac{n}{N} = \frac{Q_g/Q_o - a}{b}, \qquad a = 0.29\cos\varphi, \quad b = 0.52
\label{eq:nN}
\end{equation}
\begin{equation}
PPB_p = (107.2 + 8.604\,Q_o)\,cT_c\,\frac{n}{N} + (31.7 + 5.234\,Q_o)\,cT_n\left(1 - \frac{n}{N}\right)
\label{eq:ppbp}
\end{equation}
\begin{equation}
PP_f = \sum_{i=1}^{ciclo} \frac{PPB_{p,i}\; C_{IAF}\; C_{R,i}\; C_c}{1 - 0.01\,U}
\label{eq:ppf}
\end{equation}
where $Q_g$ and $Q_o$ are the global and extraterrestrial solar radiation (MJ m$^{-2}$ d$^{-1}$);
$\varphi$ is the latitude; $PPB_p$ is the gross potential biomass production of the standard crop
(kg DM ha$^{-1}$ d$^{-1}$); $cT_c$ and $cT_n$ are the temperature corrections for clear and overcast
skies; $C_{IAF}$ is the leaf area correction (0.5 for a maximum leaf area index of 5); $C_R$ is the
respiration correction (0.5 for $T_{med} \ge$ 20 \textdegree{}C and 0.6 below); $C_c$ is the harvest
index (0.40); and $U$ is the grain moisture (13\%).

The attainable yield ($PA$) was obtained by penalizing $PP_f$ sequentially in each yield response
period by its relative evapotranspiration deficit (eq. 10), with the yield response factors $k_y$ =
0.4, 0.4, 1.5, 0.5 and 0.2 for establishment, vegetative, flowering, grain filling and maturity
\citep{doorenbos1979, battisti2015}.
\begin{equation}
PA = PP_f \prod_{j=1}^{5} \max\!\left[0,\; 1 - k_{y,j}\left(1 - ISNA_j\right)\right]
\label{eq:pa}
\end{equation}
For each cycle, the model also produced the agroclimatic indicators used as features: $ISNA$, water
deficit, water surplus, rainfall, mean, maximum and minimum temperature, solar radiation and relative
humidity in each period, the number of days with $T_{min} \le$ 2 \textdegree{}C (frost) and with
$T_{max} \ge$ 35 \textdegree{}C (heat), the water efficiency $PA/PP_f$ and the cycle length. Frost
days were summed over flowering and grain filling (reproductive frost), and the cold severity was the
number of degrees by which the lowest $T_{min}$ of these two periods fell below 2 \textdegree{}C (zero
otherwise).

\subsection{APSIM, AquaCrop and DSSAT}
\label{subsec:otros}

Three more detailed process-based models were run without calibration, with their default parameters,
to be compared with the MZA-FAO as feature generators. All used the same BR-DWGD weather and a soil
profile whose wilting point, field capacity and saturation were estimated layer by layer from the
MapBiomas Soil clay and sand contents with the pedotransfer functions of Saxton and Rawls, with the
available water scaled to the municipal DTA \citep{saxton2006}.

APSIM Next Generation (Maize module) was compiled from the open source code (.NET 8) and run from its
official maize example, with the generic cultivar A\_100, 6 plants m$^{-2}$, 200 kg N ha$^{-1}$ at sowing
(non-limiting nitrogen), a seven-layer profile from 0 to 1.8 m and continuous simulation from 1998 to
2025, sowing every year on the scenario date. AquaCrop was run with AquaCrop-OSPy, with the default
growing-degree-day maize crop (\texttt{MaizeGDD}), $ET_o$ from BR-DWGD and 70\% of the available water at
the start of each season. DSSAT 4.8 CSM-CERES-Maize was compiled from the open source code and driven by
DSSATTools, in seasonal mode, with the generic short-season cultivar 990003, 6 plants m$^{-2}$,
200 kg N ha$^{-1}$ and soybean as the previous crop. From each model, the simulated yield and biomass,
grain number (APSIM), the ratio between simulated and potential yield (AquaCrop), cycle length, days to
anthesis (DSSAT), the cumulative water stress in the vegetative, flowering and grain filling periods
(APSIM: $\sum(1 - F_w)$; AquaCrop: $\sum(1 - T_r/T_{r,pot})$; DSSAT: $\sum WSPD$, with flowering as
$\pm$10 days around anthesis) and the minimum $T_{min}$ in each period were retained
\citep{holzworth2014, steduto2009, kelly2021, jones2003}. For a fair comparison, the four models were
compared on the same four sowing dates (15 January, 5 February, 25 February and 15 March) in the 373
municipalities with yield records.

\subsection{Hybrid yield models}
\label{subsec:hybrid}

The target was the logarithm of the municipal yield, decomposed into a technology trend (linear in the
harvest year), a municipality effect and an annual anomaly (eq. 11); the machine learning models predicted
the anomaly $\varepsilon$ from the crop model features.
\begin{equation}
\ln Y_{m,t} = \beta_0 + \beta_1\, t + \alpha_m + \varepsilon_{m,t}
\label{eq:target}
\end{equation}
where $Y_{m,t}$ is the yield of municipality $m$ in harvest year $t$ (kg ha$^{-1}$); $\beta_0$ and
$\beta_1$ are the trend coefficients; $\alpha_m$ is the municipality effect; and $\varepsilon_{m,t}$ is
the anomaly explained by weather. The trend was fitted first by ordinary least squares, and $\alpha_m$
was the mean residual of each municipality; municipalities absent from a training set received
$\alpha_m = 0$.

Since IBGE does not report sowing dates, the municipal yield mixes the sowings of the season; the
training features of each model were therefore the means of the four window dates, while predictions
for a given date used the features of that date. Eleven formulations were compared (Table~\ref{tab:models}):
the historical mean (A), a linear regression on $ISNA$ at flowering and grain filling and cold severity
(B, calibrated MZA), a gradient boosting on raw weather by period (C, pure machine learning), the same
gradient boosting on MZA-FAO features without (D) and with (E) monotonic constraints, the monotonic
gradient boosting on APSIM (F), AquaCrop (H) and DSSAT (I) features and on all of them together (J), a
support vector regression on MZA-FAO features (K) and the ensemble of E and K (L). Altitude, DTA,
latitude and longitude were included in all machine learning models. The gradient boosting was shallow
and regularized (histogram gradient boosting, maximum depth 3, learning rate 0.05, 300 iterations, at
least 100 samples per leaf, L2 penalty 1.0), because the weather signal is annual and the training set
has only 22 years. In E, F, H and I, monotonic constraints forced yield to increase with $ISNA$ and
$PA/PP_f$ (or with simulated yield and grain number) and to decrease with water deficit, water stress,
frost, cold and heat. The support vector regression used a radial basis kernel ($C$ = 0.3, $\gamma$ =
0.05, $\epsilon$ = 0.05, standardized features; $C$ and $\gamma$ were those chosen by the nested tuning
described in Section~\ref{subsec:validation} in 22 of the 23 folds), and $PA/PP_f$ and the flowering deficit were excluded
from it because they are collinear with the flowering $ISNA$ ($r$ = 0.97 and $-$0.92) and, without a
monotonic constraint, the model split the water effect between them with opposite signs. The final model
(L) was the mean of the log predictions of E and K \citep{karpatne2017, smola2004, pedregosa2011}.

% Table of hybrid model formulations
\begin{table*}[t]
\centering
\caption{Yield model formulations compared in the leave-one-year-out validation. All machine learning
models also used altitude, DTA, latitude and longitude.}
\label{tab:models}
\small
\begin{tabularx}{\textwidth}{l l X}
\toprule
Model & Algorithm & Features \\
\midrule
A & Trend + municipality effect & -- \\
B & Linear regression & $ISNA$ at flowering and grain filling, cold severity (MZA-FAO) \\
C & Gradient boosting & Rainfall, $T_{min}$ and $T_{max}$ in the vegetative, flowering and grain filling periods \\
D & Gradient boosting & $PA/PP_f$, $ISNA$ (vegetative, flowering, grain filling), flowering deficit, frost, cold and heat (MZA-FAO) \\
E & Monotonic gradient boosting & Same as D \\
F & Monotonic gradient boosting & APSIM yield, grain number, water stress per period, $T_{min}$, cycle length \\
H & Monotonic gradient boosting & AquaCrop yield, yield ratio, transpiration stress per period, $T_{min}$, cycle length \\
I & Monotonic gradient boosting & DSSAT yield, biomass, water stress per period, $T_{min}$, days to anthesis \\
J & Monotonic gradient boosting & Features of D, F, H and I \\
K & Support vector regression & Same as D, without $PA/PP_f$ and flowering deficit \\
L & Ensemble (mean of E and K) & -- \\
\bottomrule
\end{tabularx}
\end{table*}

\subsection{Validation and robustness}
\label{subsec:validation}

All models were validated by leaving one whole harvest year out (23 folds, 2003--2025): the trend, the
municipality effects and the machine learning models were estimated with the other 22 years only, which
prevents information on the weather of the test year from leaking into training. Performance was
evaluated with the coefficient of determination ($R^2$), the root mean square error ($RMSE$) and the
mean absolute error ($MAE$) of the back-transformed yields (eqs. 12--14), at the municipal scale and
after averaging observed and predicted yields by microregion, mesoregion and state in each year
\citep{roberts2017, kapoor2023}.
\begin{equation}
R^2 = 1 - \frac{\sum_{k}(O_k - E_k)^2}{\sum_{k}(O_k - \bar{O})^2}
\label{eq:r2}
\end{equation}
\begin{equation}
RMSE = \sqrt{\frac{1}{n}\sum_{k=1}^{n}(E_k - O_k)^2}
\label{eq:rmse}
\end{equation}
\begin{equation}
MAE = \frac{1}{n}\sum_{k=1}^{n}\lvert E_k - O_k \rvert
\label{eq:mae}
\end{equation}
where $O_k$ and $E_k$ are the observed and estimated yields (kg ha$^{-1}$), $\bar{O}$ is the mean
observed yield and $n$ is the number of municipality-years.

Differences in $RMSE$ between models were tested by bootstrap over years (2,000 resamples of the 23 test
years, the independent unit of the validation), reporting the median difference relative to model E and
its 95\% interval \citep{efron1993}. Two robustness analyses were performed. First, because the MZA-FAO
features were designed with agronomic reasoning, the comparison of the four feature generators was
repeated with an identical automatic selection inside each fold: the $k$ = 8 or 12 candidates with the
highest absolute Spearman correlation with the training anomaly were chosen, skipping candidates
correlated above 0.9 with one already chosen, from pools containing only process outputs of each model,
and the same unconstrained gradient boosting was used for all. Second, HGB, XGBoost, LightGBM, random
forest, extremely randomized trees, ridge regression with splines, multilayer perceptron and support
vector regression were compared on the MZA-FAO features in default, shallow and nested-tuned
configurations, the latter with the hyperparameters chosen by a grid evaluated on five blocks of years
within the training set of each outer fold, so that the choice never used the test year; the ensembles of
the best single algorithm with the monotonic gradient boosting were also evaluated.

\subsection{Prediction intervals, phase analyses and interpretability}
\label{subsec:interpret}

Prediction intervals were obtained by leave-one-year-out conformal prediction: for each test year, the
10th, 50th and 90th percentiles of the log residuals of the other years were added to the prediction,
and the coverage of the P10--P90 interval was checked year by year \citep{lei2018}. The relation between
weather and yield was analysed by phenological period at two scales, with the Pearson correlation
between the yield anomaly of model A (fitted with all years) and each variable: between years, using the
state means of each year, and between municipalities, as the mean of the correlations within each year.
The contribution of each feature to the final model was quantified with Shapley additive explanations
(SHAP); since SHAP values are additive over models, the ensemble values were the mean of the exact tree
SHAP values of the gradient boosting and of the permutation SHAP values of the support vector regression,
and they were grouped by phenological period \citep{lundberg2020}. Because the support vector regression
has no monotonic constraint, the physical consistency of the ensemble was checked with partial dependence
curves of each constrained feature between its 5th and 95th percentiles.

The final model, trained with all years, predicted the yield quantiles of each municipality, year
(1999--2025) and sowing date at the technology level of 2025 (trend evaluated in 2025), so that the 27
weather years became comparable scenarios for the risk analysis. The 26 municipalities without
second-season yield records, located on the cold southern plateau, were excluded from the decision
analysis.

\subsection{Production costs and prices}
\label{subsec:costs}

Production costs of second-season maize from 2000 to 2025 were obtained from the historical series of
the Department of Rural Economics of Paraná (SEAB/DERAL), reported per 60-kg bag and converted to
R\$ ha$^{-1}$ with the reference yield of each period. The operating cost, used in the main scenario,
was the total cost minus the factor income (imputed remuneration of land and own capital), whose share
of the total cost was taken from the CONAB cost series for Paraná in each year (linearly interpolated
in the years without CONAB data), because the decision to
sow the second season after soybean does not pay for the land again; the total cost was used as a
sensitivity scenario. The itemized DERAL cost showed that only external transport is proportional to
production, so each cost was split into a part per hectare, committed at sowing, and a part per kilogram.
Monthly prices received by farmers in Paraná (SEAB/DERAL, via Ipeadata) were used. Costs and prices were
deflated to December 2025 with the monthly IPCA, and every series was indexed by the harvest year.
[INSERT: references for the DERAL, CONAB, Ipeadata and IPCA series.]

\subsection{Field-level yield risk and margin simulation}
\label{subsec:margin}

The yield distribution of each municipality, year and sowing date was assumed lognormal, with the median
given by the model and the dispersion $\sigma$ derived from the P10--P90 interval. Because the model
predicts the municipal mean, whereas each field varies more, an idiosyncratic field-level deviation
$\tau$ was added (eq. 15). $\tau$ was the only calibrated parameter: it minimized the squared difference
between the annual frequency of field yields below 65\% of the expected yield (ZARC-suitable dates of
the 15 January--15 March window) and the fraction of second-season maize policies indemnified by the
federal crop insurance subsidy programme (PSR) in Paraná from 2016 to 2024, searched over a grid from 0
to 0.8 in steps of 0.01. The correlation between the
two annual series was used as an independent check, because $\tau$ only adjusts their mean level.
\begin{equation}
\ln Y^{f} = \ln P50 + \sqrt{\sigma^2 + \tau^2}\;\xi, \qquad
\sigma = \frac{\ln P90 - \ln P10}{2\,z_{0.9}}, \qquad \xi \sim N(0,1)
\label{eq:field}
\end{equation}
where $Y^{f}$ is the field yield (kg ha$^{-1}$); $P10$, $P50$ and $P90$ are the predicted yield
quantiles; $z_{0.9}$ is the 0.9 quantile of the standard normal distribution; and $\tau$ is the
field-level standard deviation in log.

For each municipality, sowing date and harvest year (2000--2025), 200 field yields were drawn and the
margin was computed with the price and cost of the same year (eq. 16), so that the dependence between
weather, yield, price and cost of each season was preserved. The price was that of the month of sale,
20 days after the simulated physiological maturity (time for field drying from about 30\% to 18--22\%
grain moisture). Only ZARC-suitable combinations had access to insurance: yield insurance with coverage of
65\% of the expected yield (the level of 75\% of the second-season maize policies subsidized in Paraná
from 2016 to 2025), with the indemnity of eq. 17 and the premium included in the cost; in unsuitable
combinations there was no indemnity and no premium.
\begin{equation}
M = p\,Y^{f}(1 - F) - \left(C_{ha} + c_{kg}\,Y^{f}\right) + I
\label{eq:margin}
\end{equation}
\begin{equation}
I = \max\!\left(0,\; 0.65\,\bar{Y} - Y^{f}\right)\bar{p}
\label{eq:indemnity}
\end{equation}
where $M$ is the margin (R\$ ha$^{-1}$); $p$ is the real price of the month of sale (R\$ kg$^{-1}$); $F$ is
the Funrural levy on gross revenue (1.5\%); $C_{ha}$ and $c_{kg}$ are the per-hectare (R\$ ha$^{-1}$) and
per-kilogram (R\$ kg$^{-1}$) costs; $I$ is the insurance indemnity (R\$ ha$^{-1}$); $\bar{Y}$ is the expected
yield of the municipality and sowing date (mean $P50$ over years, kg ha$^{-1}$); and $\bar{p}$ is the mean
price of the month of sale (R\$ kg$^{-1}$).

\subsection{Official risk zoning and decision criteria}
\label{subsec:decision}

The official ZARC risk table for rainfed second-season maize in Paraná, 2026/2027 season, was obtained
from the open data of the Ministry of Agriculture: for each municipality and soil class, a sowing date was
suitable when its ten-day period had an indicated risk of 20, 30 or 40\%. Cycle Group I (earliest cycle)
was used in the main scenario and Group II as sensitivity \citep{monteiro2021, embrapa2023}.
[INSERT: references for the ZARC ordinance and open data.] As a test of the MZA-FAO water balance, the
water criterion of the zoning was also replicated (100-day reference cycle; adverse year when $ISNA$ was
$\le$ 0.70 in phase I or $\le$ 0.55 in phase III; suitable date when at most 40\% of the years were
adverse) and compared with the official table.

The sowing dates were ranked by the certainty equivalent ($CE$) of the margin with a negative exponential
utility (constant absolute risk aversion), over relative risk aversion coefficients $r_r$ from 0 to 4,
following the stochastic efficiency with respect to a function (SERF) (eq. 18). The absolute coefficient
was $r_a = r_r/W$, with $W$ the mean total cost (R\$ 4,239 ha$^{-1}$) as the capital at risk, and the best
date of each municipality was the ZARC-suitable date with the highest $CE$ at $r_r$ = 2 (rather risk
averse). The probability of loss, $P(M < 0)$, and the conditional value at risk (eq. 19) complemented the
analysis, and the value of insurance under the ZARC was the difference between the $CE$ with and without
insurance \citep{hardaker2004, anderson1992, rockafellar2000}.
\begin{equation}
CE = -\frac{1}{r_a}\ln\!\left[\frac{1}{N}\sum_{s=1}^{N} \exp\left(-r_a M_s\right)\right]
\label{eq:ce}
\end{equation}
\begin{equation}
CVaR_{10} = \frac{1}{\lceil 0.1N \rceil}\sum_{s=1}^{\lceil 0.1N \rceil} M_{(s)}
\label{eq:cvar}
\end{equation}
where $M_s$ is the margin of draw $s$ (R\$ ha$^{-1}$); $N$ is the number of draws (26 years $\times$ 200);
$r_a$ is the absolute risk aversion coefficient (ha R\$$^{-1}$); and $M_{(s)}$ are the margins in
ascending order.

The ability of the yield models to anticipate economic losses was evaluated out of sample with the
leave-one-year-out predictions: the predicted probability of loss of each municipality-year (fraction of
the residuals of the other years below the break-even yield) was compared with the observed loss (IBGE
yield below the break-even yield, with the price of the typical harvest month and the operating cost of
the year), by the Brier score (eq. 20) and reliability curves \citep{brier1950}.
\begin{equation}
BS = \frac{1}{n}\sum_{k=1}^{n}\left(\hat{q}_k - o_k\right)^2
\label{eq:brier}
\end{equation}
where $\hat{q}_k$ is the predicted probability of loss and $o_k$ equals 1 when a loss was observed and 0
otherwise.

\subsection{Economic value of the yield models}
\label{subsec:regret}

The choice of the sowing date cannot be verified against observations, because the municipal yield
statistics do not report sowing dates. The decision that can be verified is whether to sow second-season
maize in each municipality and year, which was used to measure the economic value of each yield model.
With the leave-one-year-out predictions of each model, the farmer was assumed to sow when the expected
margin was positive, with the expected yield given by the model prediction combined with the residuals of
the other years, the expected price equal to the mean real price of the harvest month in the other years
and the operating cost of the year, known at sowing. The outcome was the realized margin with the observed
yield and price, and the regret of each model was the mean difference to the best decision in hindsight
(eq. 21), compared with the strategy of always sowing and with perfect information; differences between
models were tested by bootstrap over years. Because the models use the weather of the season itself, this
analysis measures the economic value of the information carried by each model rather than the skill of
a forecast issued at sowing.
\begin{equation}
R = \frac{1}{n}\sum_{k=1}^{n}\left[\max(0, M_k) - \delta_k M_k\right], \qquad
\delta_k =
\begin{cases}
1, & \widehat{\mathbb{E}}[M_k] > 0\\
0, & \text{otherwise}
\end{cases}
\label{eq:regret}
\end{equation}
where $R$ is the regret (R\$ ha$^{-1}$); $M_k$ is the realized margin of municipality-year $k$, computed
with eq. 16 using the observed yield and without insurance (R\$ ha$^{-1}$); $\delta_k$ is the decision to sow;
and $\widehat{\mathbb{E}}[M_k]$ is the expected margin at sowing.

\subsection{Sensitivity analysis}
\label{subsec:sensitivity}

The robustness of the recommendation was assessed by rerunning the chain from the step where each
parameter enters: thermal time of 1,350 and 1,550 \textdegree{}C d (MZA-FAO, hybrid model and decision);
training features weighted towards early sowing (weights 0.4, 0.3, 0.2 and 0.1 for 15 January,
5 February, 25 February and 15 March) or averaged over the seven dates of the ZARC window (hybrid model and
decision); and, in the economic step only, total instead of operating cost, insurance coverage of 70\%,
ZARC cycle Group II, 10 or 30 days from maturity to sale, and a risk-rated premium. In the main scenario
the premium was the DERAL item (a fixed share of the variable cost); in the risk-rated scenario it was
actuarially fair, equal to the expected indemnity of each municipality and sowing date (mean of eq. 17
over years and draws). For each scenario, the agreement of the best date with the main scenario was
computed over the municipalities.


% =========================================================
% FIGURE — TWO COLUMNS
% =========================================================
% Example:
% \begin{figure*}[t]
% \centering
% \includegraphics[width=0.96\textwidth]{figure1.png}
% \caption{Insert figure caption here.}
% \label{fig:figure1}
% \end{figure*}

% =========================================================
% FIGURE — ONE COLUMN
% =========================================================
% Example:
% \begin{figure}[t]
% \centering
% \includegraphics[width=\linewidth]{figure2.png}
% \caption{Insert figure caption here.}
% \label{fig:figure2}
% \end{figure}

% =========================================================
% TABLE — TWO COLUMNS
% =========================================================
% Example:
% \begin{table*}[t]
% \centering
% \caption{Insert table caption here.}
% \label{tab:table1}
% \begin{tabular}{llll}
% \toprule
% Column 1 & Column 2 & Column 3 & Column 4 \\
% \midrule
% -- & -- & -- & -- \\
% \bottomrule
% \end{tabular}
% \end{table*}


% =========================================================
% 3 RESULTS
% =========================================================
\section{Results}
\label{sec:results}

\subsection{Growing conditions simulated by sowing date}
\label{subsec:res_sim}

Delaying sowing lengthened the simulated cycle and moved the reproductive phases into autumn and winter
(Figure~\ref{fig:fases}). The mean cycle increased from 105 days for sowing on 5 January to 158 days on
25 March (+51\%), because the thermal time accumulated more slowly as temperature declined. The minimum
$T_{min}$ during grain filling fell from 14.8 to 3.6 \textdegree{}C, and the number of frost days in the
reproductive phase increased from 0.04 to 1.55. The flowering $ISNA$ decreased from 0.73 to 0.58 between 5
January and 25 February, when flowering coincided with the lowest rainfall of the season, and recovered
partially for later dates, which flowered in the cooler and less demanding conditions of May. As a result,
the attainable yield simulated by the MZA-FAO fell from 2,933 kg ha$^{-1}$ on 5 January to 1,738 kg
ha$^{-1}$ on 25 February ($-$41\%), whereas the potential yield decreased by only 8\% (from 6,932 to
6,349 kg ha$^{-1}$); the water efficiency $PA/PP_f$ decreased from 0.43 to 0.28.

\begin{figure*}[t]
\centering
\includegraphics[width=0.96\textwidth]{figures/fig_variaveis_fase_data.png}
\caption{Mean MZA-FAO variables by phenological period and sowing date for the 399 municipalities of
Paraná, 1999--2025: water requirement satisfaction index ($ISNA$), water deficit, minimum $T_{min}$, days
with $T_{min} \le$ 2 \textdegree{}C, rainfall and global solar radiation. Estab., establishment; Veg.,
vegetative; Flow., flowering; Fill., grain filling; Mat., maturity.}
\label{fig:fases}
\end{figure*}

\subsection{Performance of the yield models}
\label{subsec:res_models}

The historical mean (model A) explained 26\% of the variance of the municipal yields and 31\% of the
annual variation of the state mean (Table~\ref{tab:validation}), which corresponds to the part of the yield
attributable to the technology trend and to persistent differences between municipalities. All
weather-based models improved on it, and the hybrid models with MZA-FAO features were the most accurate.
With the same algorithm and number of features, the MZA-FAO features (model D) raised the $R^2$ from 0.318
to 0.403 and reduced the $RMSE$ from 1,158 to 1,084 kg ha$^{-1}$ ($-$6.4\%) in relation to raw weather
by period (model C). The monotonic constraints (model E) cost 0.009 in $R^2$ relative to D, a difference
within the uncertainty of the bootstrap ($-$7 kg ha$^{-1}$; 95\% interval $-$44 to 27). The final ensemble
(L) reached $R^2$ = 0.427, $RMSE$ = 1,062 kg ha$^{-1}$ and $MAE$ = 826 kg ha$^{-1}$, with an $RMSE$ 12\%
lower than that of the historical mean, and was the only formulation whose improvement over E excluded zero
($-$30 kg ha$^{-1}$; $-$52 to $-$9). Performance increased with spatial aggregation: the $R^2$ of the final
model was 0.540 for microregions, 0.611 for mesoregions and 0.674 for the state mean between years.

\begin{table*}[t]
\centering
\caption{Leave-one-year-out performance of the yield models (6,806 municipality-years, 2003--2025).
$\Delta RMSE$, median difference in $RMSE$ relative to model E, with the 95\% bootstrap interval over
years. Model codes as in Table~\ref{tab:models}.}
\label{tab:validation}
\small
\begin{tabular}{lrrrrrrl}
\toprule
Model & $R^2$ & $R^2$ micro. & $R^2$ meso. & $R^2$ state & $RMSE$ (kg ha$^{-1}$) & $MAE$ (kg ha$^{-1}$) & $\Delta RMSE$ vs E \\
\midrule
A Historical mean   & 0.259 & 0.385 & 0.400 & 0.313 & 1,207 & 956 & $+$108 ($-$57; 285) \\
B Calibrated MZA    & 0.353 & 0.474 & 0.547 & 0.587 & 1,128 & 883 & $+$37 ($-$36; 103) \\
C Pure ML           & 0.318 & 0.465 & 0.511 & 0.554 & 1,158 & 910 & $+$65 ($-$23; 153) \\
D MZA hybrid        & 0.403 & 0.530 & 0.598 & 0.656 & 1,084 & 848 & $-$7 ($-$44; 27) \\
E MZA hybrid, monotonic & 0.394 & 0.514 & 0.584 & 0.634 & 1,092 & 850 & -- \\
F APSIM hybrid      & 0.334 & 0.454 & 0.487 & 0.513 & 1,145 & 905 & $+$51 ($-$1; 110) \\
H AquaCrop hybrid   & 0.245 & 0.420 & 0.485 & 0.422 & 1,218 & 958 & $+$123 (0; 254) \\
I DSSAT hybrid      & 0.334 & 0.455 & 0.492 & 0.523 & 1,144 & 895 & $+$52 (11; 94) \\
J All models        & 0.350 & 0.480 & 0.532 & 0.560 & 1,130 & 887 & $+$39 (10; 65) \\
K MZA hybrid, SVR   & 0.420 & 0.531 & 0.603 & 0.682 & 1,068 & 831 & $-$24 ($-$70; 19) \\
L Ensemble E + K    & 0.427 & 0.540 & 0.611 & 0.674 & 1,062 & 826 & $-$30 ($-$52; $-$9) \\
\bottomrule
\end{tabular}
\end{table*}

The uncalibrated crop models did not improve the hybrid model. With APSIM and DSSAT features (models F and
I), the $R^2$ was 0.334, close to that of the calibrated MZA (0.353) and 51--52 kg ha$^{-1}$ worse in $RMSE$
than with MZA-FAO features; with AquaCrop (H), the hybrid model was not better than the historical mean
($R^2$ = 0.245). In AquaCrop, the default temperate maize crop failed in about half of the simulated
combinations because of cold stress at pollination. Combining the features of the four models (J) reduced
the accuracy ($R^2$ = 0.350; $+$39 kg ha$^{-1}$, 10 to 65), because more features had to be learned from
the same 22 training years. When the same automatic feature selection was applied to all generators, the
differences between MZA-FAO, APSIM and DSSAT disappeared ($R^2$ of 0.326, 0.319 and 0.330 with eight
features; $\Delta RMSE$ relative to MZA of $+$6 and $-$4 kg ha$^{-1}$, with intervals including zero), and
AquaCrop remained the worst (0.245). Among algorithms trained on the MZA-FAO features, the $R^2$ ranged from
0.294 to 0.437 (Supplementary Table S1): library default configurations were the least accurate (0.294 for
XGBoost and 0.366--0.377 for random forest, extremely randomized trees, LightGBM and HGB), nested tuning
changed the shallow gradient boosting by about 0.01, and the support vector regression was the best single
algorithm (0.436 with nested tuning).

The annual means of the final model followed the observed state yield in most years (Figure~\ref{fig:anual}),
but the largest errors were concentrated in 2015 ($-$852 kg ha$^{-1}$), 2020 ($-$1,286), 2021 ($+$841)
and 2022 ($+$907). The 80\% conformal interval covered 79.4\% of the observations overall, with coverage
between 0.62 and 0.91 in most years, but only 0.41, 0.54 and 0.69 in 2020, 2021 and 2022.

\begin{figure*}[t]
\centering
\includegraphics[width=0.96\textwidth]{figures/fig_validacao_anual.png}
\caption{Observed (IBGE, mean of the municipalities) and leave-one-year-out predicted second-season maize
yield in Paraná, 2003--2025, for the eleven model formulations.}
\label{fig:anual}
\end{figure*}

\subsection{Weather and yield by phenological phase}
\label{subsec:res_phases}

The yield anomaly was strongly related to the water balance between years, but not between municipalities
within a year (Figure~\ref{fig:cor}). Between years, the correlation with $ISNA$ was highest at flowering
($r$ = 0.73), followed by grain filling (0.66) and the vegetative period (0.65), and the water deficit showed
the symmetric pattern ($r$ = $-$0.66 at flowering). Frost became more harmful as the crop advanced
($r$ = $-$0.05 in the vegetative period, $-$0.39 at grain filling and $-$0.48 at maturity), and solar
radiation and maximum temperature were negatively correlated with yield ($r$ from $-$0.34 to $-$0.42 for
radiation), because sunny and hot seasons were also dry seasons. Between municipalities, no correlation
exceeded 0.11 in absolute value. In the final model, the static variables (altitude, DTA and location)
accounted for 28\% of the mean absolute SHAP value, flowering for 27\%, grain filling for 21\%, the
reproductive frost and cold indicators for 16\% and the vegetative period for 8\% (Figure~\ref{fig:shap}).
The most important single feature was the $ISNA$ at grain filling, followed by reproductive frost, $ISNA$
at flowering and heat at flowering.

All constrained variables kept their physical direction in the ensemble: across the 5th--95th percentile
range, the partial dependence increased with $ISNA$ (total effect of 0.100, 0.134 and 0.244 in log for the
vegetative, flowering and grain filling periods) and decreased with frost (0.159), heat (0.128), cold
(0.046) and flowering deficit (0.045). The only deviation was a flattening of the response to the grain
filling $ISNA$ at its highest values. The ensemble and the monotonic model E indicated the same sowing date
of highest mean yield in 86\% of the municipalities.

\begin{figure*}[t]
\centering
\includegraphics[width=0.96\textwidth]{figures/fig_correlacao_fase.png}
\caption{Pearson correlation between the yield anomaly (residual of the trend and municipality effect) and
the MZA-FAO variables of each phenological period, between years (state means) and between municipalities
(mean of the within-year correlations), 2003--2025.}
\label{fig:cor}
\end{figure*}

\begin{figure*}[t]
\centering
\includegraphics[width=0.96\textwidth]{figures/fig_shap_fase.png}
\caption{Contribution of each phenological period (left) and of each feature (right) to the final hybrid
model, as the mean absolute SHAP value of the ensemble of the monotonic gradient boosting and the support
vector regression.}
\label{fig:shap}
\end{figure*}

\subsection{Yield, economic risk and the best sowing date}
\label{subsec:res_risk}

At the 2025 technology level, the median predicted yield decreased from 5,691 kg ha$^{-1}$ for sowing on
5 January to 5,050 kg ha$^{-1}$ on 25 February ($-$11\%) and increased slightly for later dates (5,309 kg
ha$^{-1}$ on 25 March). In the official ZARC, 395 municipalities had 5 January as a suitable date (333 with
20\% risk), against 274 on 15 February and 226 on 15 March, and no municipality could sow on 25 March. The
replica of the ZARC water criterion with the MZA-FAO agreed with the official classification in only 39\%
of the municipality-date combinations. The real price of the month of sale decreased from R\$ 0.84 kg$^{-1}$
for the earliest sowing, sold in April--May, to R\$ 0.81 kg$^{-1}$ for sowings sold in July--August, in
line with the seasonal price index (1.006 in April and 0.952 in July; Supplementary Figure S1).

Early sowing combined the highest mean margin with the lowest risk (Table~\ref{tab:risk}). The mean
margin over operating cost decreased from R\$ 1,806 ha$^{-1}$ on 5 January to R\$ 1,083 ha$^{-1}$ on
25 February, the probability of loss increased from 0.26 to 0.37, the CVaR$_{10}$ fell from $-$R\$ 999
to $-$R\$ 1,466 ha$^{-1}$, and the certainty equivalent at $r_r$ = 2 was 2.5 times higher on 5 January
(R\$ 853 ha$^{-1}$) than on 25 February (R\$ 338 ha$^{-1}$). Over the total cost, the probability of loss
was 0.46 on 5 January and 0.55--0.57 from February onwards. Insurance access under the ZARC was worth
R\$ 39 to R\$ 75 ha$^{-1}$ in certainty equivalent, with the highest value for the earliest dates.

\begin{table*}[t]
\centering
\caption{Mean economic risk indicators by sowing date over the 372 municipalities with second-season
records (operating cost; 26 years $\times$ 200 field yields). $CE$, certainty equivalent at relative risk
aversion 2; insurance value, $CE$ with minus $CE$ without insurance (ZARC-suitable dates only).}
\label{tab:risk}
\small
\begin{tabular}{lrrrrrr}
\toprule
Sowing date & Margin (R\$ ha$^{-1}$) & $P$(loss) & CVaR$_{10}$ (R\$ ha$^{-1}$) & $CE$ (R\$ ha$^{-1}$) & Insurance value (R\$ ha$^{-1}$) & Best date (n) \\
\midrule
5 Jan  & 1,806 & 0.26 & $-$999   & 853 & 75 & 311 \\
15 Jan & 1,628 & 0.28 & $-$1,108 & 738 & 59 & 31 \\
25 Jan & 1,420 & 0.30 & $-$1,229 & 612 & 50 & 1 \\
5 Feb  & 1,257 & 0.34 & $-$1,355 & 465 & 64 & 2 \\
15 Feb & 1,133 & 0.36 & $-$1,403 & 377 & 58 & 0 \\
25 Feb & 1,083 & 0.37 & $-$1,466 & 338 & 50 & 0 \\
5 Mar  & 1,176 & 0.36 & $-$1,503 & 388 & 53 & 0 \\
15 Mar & 1,319 & 0.33 & $-$1,487 & 472 & 39 & 27 \\
25 Mar & 1,468 & 0.31 & $-$1,544 & 527 & -- & 0 \\
\bottomrule
\end{tabular}
\end{table*}

Restricted to ZARC-suitable dates, the best date was 5 January in 311 of the 372 municipalities (84\%),
15 January in 31 and 15 March in 27, while 27 municipalities of the southern plateau had no suitable date
(Figure~\ref{fig:mapas}). The mean certainty equivalent at the best date was R\$ 866 ha$^{-1}$, with a
probability of loss of 0.26. Without the ZARC restriction, 40 municipalities would choose 25 March, and
the best date changed between risk neutrality and $r_r$ = 2 in 86 municipalities (Figure~\ref{fig:serf}).

The 27 municipalities that preferred 15 March were all located on the western border along the Paraná
River and the Itaipu reservoir and in the sandy northwest, at a mean altitude of 299 m, against 583 m in
the municipalities that preferred 5 January. In these municipalities, sowing on 5 January led to flowering
under a mean $T_{max}$ of 32.3 \textdegree{}C, with 1.9 days above 35 \textdegree{}C, against 27.2
\textdegree{}C and 0.2 days for sowing on 15 March, whereas reproductive frost remained negligible (0.29
days). The median predicted yield was 9\% higher on 15 March, and replacing one feature at a time by its
15 March value attributed 0.076 of the 0.087 difference in log yield to heat at flowering, against 0.017
for the grain filling $ISNA$ and $-$0.010 for frost. The preference did not come from prices: with the same
price for every date, the advantage of 15 March in certainty equivalent increased from R\$ 139 to
R\$ 323 ha$^{-1}$. The MZA-FAO alone would have favoured early sowing in these municipalities (attainable
yield of 1,932 kg ha$^{-1}$ on 5 January against 1,339 kg ha$^{-1}$ on 15 March).

\begin{figure*}[t]
\centering
\includegraphics[width=0.96\textwidth]{figures/fig_mapas_risco_decisao.png}
\caption{(a) Official ZARC risk class for sowing on 15 February; (b) best sowing date by certainty
equivalent ($r_r$ = 2) among ZARC-suitable dates; (c) certainty equivalent at the best date; (d) value of
insurance under the ZARC for sowing on 15 February. Grey, municipalities without second-season records or
without suitable dates.}
\label{fig:mapas}
\end{figure*}

\begin{figure*}[t]
\centering
\includegraphics[width=0.96\textwidth]{figures/fig_serf_exemplos.png}
\caption{Stochastic efficiency with respect to a function (SERF): certainty equivalent of each sowing
date as a function of the relative risk aversion in four municipalities. Dotted lines, dates outside the
ZARC.}
\label{fig:serf}
\end{figure*}

\subsection{Independent checks of the risk estimates}
\label{subsec:res_checks}

The calibrated field-level deviation was $\tau$ = 0.33 in log, larger than the mean municipal dispersion
($\sigma$ = 0.26). With it, the modelled frequency of field yields below 65\% of the expected yield
followed the fraction of indemnified PSR policies between years ($r$ = 0.88; Figure~\ref{fig:psr}),
including the high losses of 2018 and 2021 and the low losses of 2017 and 2023, although the model
underestimated the extreme of 2021 (0.47 against 0.68). Out of sample, 55\% of the municipality-years had
a yield below the break-even yield over operating cost. All models predicted this event better than a
constant frequency (Brier score of 0.248): the hybrid models with MZA-FAO features scored 0.164--0.171 (0.165
for the final model), the historical mean 0.173 and the AquaCrop hybrid 0.180 (Supplementary Figure S2).

\begin{figure}[t]
\centering
\includegraphics[width=\linewidth]{figures/fig_validacao_psr.png}
\caption{Annual frequency of yields below 65\% of the expected yield at the municipal and field levels
(ZARC-suitable dates, 15 January--15 March) and fraction of indemnified second-season maize policies of
the federal insurance subsidy programme (PSR) in Paraná, 2016--2024.}
\label{fig:psr}
\end{figure}

In the plant/do-not-plant decision evaluated with observed yields, always sowing resulted in a mean
regret of R\$ 448 ha$^{-1}$ and a realized margin close to zero ($-$R\$ 4 ha$^{-1}$), and perfect
information would have yielded R\$ 444 ha$^{-1}$. Deciding with any of the yield models reduced the regret
to R\$ 307--338 ha$^{-1}$ (58--60\% of correct decisions, against 45\% when always sowing), although the
bootstrap intervals of this reduction included zero (e.g. $-$R\$ 121 ha$^{-1}$, $-$317 to 71, for the
final model). The models did not differ among themselves: the historical mean had the lowest regret
(R\$ 307 ha$^{-1}$), and the difference between the final model and model E was $-$R\$ 3 ha$^{-1}$
($-$19 to 21).

\subsection{Sensitivity of the recommendation}
\label{subsec:res_sens}

The recommendation of early sowing was robust to all scenarios (Table~\ref{tab:sens},
Figure~\ref{fig:sens}). The economic scenarios kept the same best date in 90--99\% of the municipalities:
the total cost reduced the mean margin to R\$ 405 ha$^{-1}$ and raised the probability of loss to 0.52,
but changed the best date in only 2\% of the municipalities. A risk-rated premium would cost R\$ 123
ha$^{-1}$ on average (R\$ 94--154 ha$^{-1}$ between the 10th and 90th percentiles), less than the fixed
DERAL premium (R\$ 169 ha$^{-1}$), raising the value of insurance under the ZARC to R\$ 101--129 ha$^{-1}$
with the same best date in 98\% of the municipalities. The scenarios that changed the yield model had
larger effects on the best date (78--91\% agreement), but 5 January remained the best date in 284 to 323
municipalities. A thermal time of 1,550 \textdegree{}C d improved the hybrid model ($R^2$ = 0.450 and annual
$R^2$ = 0.727), whereas 1,350 \textdegree{}C d and training features weighted towards early sowing
reduced its accuracy ($R^2$ of 0.393 and 0.384).

\begin{table}[t]
\centering
\caption{Sensitivity of the hybrid model and of the best sowing date. Agreement, share of municipalities
with the same best date as the main scenario.}
\label{tab:sens}
\small
\begin{tabular}{lrrr}
\toprule
Scenario & $R^2$ & Agreement (\%) & Best 5 Jan / 15 Mar (n) \\
\midrule
Main scenario                 & 0.427 & --   & 311 / 27 \\
Thermal time 1,350 \textdegree{}C d & 0.393 & 78 & 284 / 12 \\
Thermal time 1,550 \textdegree{}C d & 0.450 & 83 & 310 / 23 \\
Early-weighted window         & 0.384 & 91 & 310 / 24 \\
Seven-date ZARC window        & 0.430 & 90 & 323 / 26 \\
Total cost                    & --    & 98 & -- \\
Coverage 70\%                 & --    & 99 & -- \\
ZARC cycle Group II           & --    & 95 & -- \\
10 / 30 days to sale          & --    & 90 & -- \\
Risk-rated premium            & --    & 98 & -- \\
\bottomrule
\end{tabular}
\end{table}

\begin{figure}[t]
\centering
\includegraphics[width=\linewidth]{figures/fig_sensibilidade_ce.png}
\caption{Mean certainty equivalent ($r_r$ = 2) by sowing date in the main scenario and in the scenarios
that change the crop model or the training of the hybrid model.}
\label{fig:sens}
\end{figure}


% =========================================================
% 4 DISCUSSION
% =========================================================
\section{Discussion}
\label{sec:discussion}

\subsection{What limits the prediction of municipal yields}
\label{subsec:dis_limits}

The final hybrid model explained 43\% of the variance of municipal yields and 67\% of the annual variation
of the state mean, values that depend strongly on the scale. Most of the remaining error is variation
between municipalities within the same year, which the weather variables did not explain: between
municipalities, no correlation between the yield anomaly and the water balance exceeded 0.11, whereas
between years it reached 0.73. Municipal yields from survey statistics combine management, cultivar choice,
pests and estimation error, and repeated values (8.5\% of the records) show that part of them are estimates
rather than measurements; this noise is reduced by aggregation, which explains the increase of the $R^2$
from 0.43 in municipalities to 0.61 in mesoregions. Values of this order are consistent with the third of
the yield variability attributed to climate worldwide, and the gap between municipal and aggregated
accuracy agrees with the view that statistical yield models capture the large-scale weather signal better
than local variation \citep{ray2015, lobell2010}.

The validation design is as important as the algorithm for interpreting these values. When
municipality-years of the same season are split at random, the model sees the weather of the test year
in other municipalities and the accuracy is inflated; leaving whole years out, as here, measures the
ability to predict a season not seen in training, which is what matters for risk analysis
\citep{roberts2017, kapoor2023}. The largest errors were concentrated in a few seasons. In 2021, all models
overestimated the yield, and the 80\% interval covered only 54\% of the observations, as in a season in
which drought and frost acted together and the PSR recorded indemnities in 68\% of the policies. In 2022,
the overestimation occurred without an equivalent weather signal, which suggests a non-climatic cause; the
outbreak of corn leafhopper and corn stunt diseases reported for the 2021/22 second season in southern
Brazil is a plausible explanation that the model cannot represent. In 2020, conversely, the models predicted
a drought loss that did not appear in the IBGE statistics, which indicates that part of the error may lie
in the official estimates themselves.

\subsection{Simple and complex crop models as feature generators}
\label{subsec:dis_models}

The hybrid model did not benefit from more detailed crop models used without calibration. APSIM and DSSAT
produced hybrid models as accurate as the calibrated MZA-FAO, AquaCrop was worse than the historical mean,
and with an identical automatic feature selection, MZA-FAO, APSIM and DSSAT performed alike. This result
can be explained by three factors. First, the default cultivars of the detailed models were developed for
other environments: the AquaCrop maize crop is parameterized for temperate summer maize and failed in about
half of the second-season cycles because of cold stress at pollination, and the generic DSSAT cultivar
simulated cycles of 85--115 days with yield almost insensitive to the sowing date. Second, the
additional processes of these models (nitrogen, canopy development, kernel set) were either non-limiting
in the simulations or not informed by local data, so they added variance rather than information. Third,
the target is an annual anomaly learned from only 22 seasons, which penalizes any increase in the number of
features, as shown by the loss of accuracy when the four models were combined.

These results agree with model intercomparisons that found large divergence among uncalibrated crop models
for maize and among soybean models in southern Brazil, and with the conclusion that crop model performance
depends more on calibration than on structure \citep{bassu2014, battisti2017, seidel2018}. The advantage of
the hybrid MZA-FAO in the main comparison came from agronomically designed features (water satisfaction
by period, water efficiency and reproductive frost, cold and heat indicators) and from the physical
constraints; since these features were refined on the same data, the robustness analysis with automatic
selection is the fairer comparison of the generators. Previous hybrids improved yield prediction with
calibrated APSIM in Australia and in the US Corn Belt \citep{feng2019, shahhosseini2021}; the present
results show that, where calibration data are lacking, a simple model with few parameters delivers the
same information at a much lower cost, which favours its operational use in yield monitoring and risk
zoning \citep{battisti2015, noia2024}.

\subsection{Physical guidance and interpretability}
\label{subsec:dis_physics}

The relations learned by the hybrid model are agronomically coherent. The flowering period had the
highest correlation with the yield anomaly and, together with grain filling, concentrated about two thirds
of the SHAP contribution, which matches the high sensitivity of maize to water deficit at flowering
($k_y$ = 1.5) and the importance of grain filling for second-season maize that completes its cycle under
declining rainfall \citep{doorenbos1979, andrea2018}. Drought during grain filling and heat before anthesis
were also identified as the most damaging extremes for wheat with a hybrid APSIM model \citep{feng2019}.
The monotonic constraints ensured that the response to each variable had the physical direction expected
when the model was applied to individual sowing dates, a condition that purely statistical models do not
guarantee \citep{karpatne2017}. The support vector regression added accuracy but required removing
variables collinear with the flowering $ISNA$, because without constraints it split the water effect with
opposite signs; averaging it with the constrained model kept the physical consistency of the response.

\subsection{Early sowing and the exception of the warm lowlands}
\label{subsec:dis_sowing}

The economic analysis confirmed early sowing as the best strategy for most of Paraná. Each delay pushed
flowering and grain filling into drier and colder conditions, reducing the mean margin, increasing the
probability of loss and, above all, lowering the certainty equivalent, which was 2.5 times higher on
5 January than on 25 February. This agrees with the risk zoning of second-season maize in Paraná, which
identified frost during grain filling as the main limitation in the south-central and western regions,
and with the conclusion that anticipating the soybean sowing is the best strategy to improve the total
production of the soybean--maize succession \citep{goncalves2002, noia2019eja, noia2020}. In practice, sowing
on 5 January requires very early soybean cultivars sown in September, whose yield and costs were not
modelled here; the recommendation should therefore be read as ``as early as the soybean harvest allows''
rather than as a fixed calendar date.

The western border and the sandy northwest were the exception, preferring sowing in mid-March. The
attribution analysis showed that this preference came from heat at flowering: in these warm lowlands, early
sowing makes flowering coincide with the hottest weeks of the season, whereas late sowing flowers under
milder temperatures without exposure to frost. Maize yield declines sharply when temperatures exceed about
30 \textdegree{}C during the season, and heat around flowering impairs pollination and kernel set
\citep{schlenker2009, lobell2011}. The MZA-FAO alone does not penalize heat, so it favoured early sowing in
these municipalities; the heat effect was learned by the machine learning from the variation between years,
which is a concrete example of what the hybrid adds to the physics. Because this effect is extrapolated from
the window mean to individual dates, and late sowing still carries some frost risk and harvest in winter, the
recommendation for these municipalities should be confirmed with field trials.

\subsection{Economic risk, insurance and the value of the yield models}
\label{subsec:dis_economics}

Accounting for risk changed the decision in a meaningful part of the state: the best date differed between
risk neutrality and moderate risk aversion in 86 municipalities, and without the ZARC restriction 40
municipalities would choose 25 March, a date with high mean margin in some years but a heavy lower tail.
Ranking the dates by the certainty equivalent over a range of risk aversion therefore gives a more complete
basis for advice than the mean yield or the mean profitability of each date \citep{hardaker2004, oligini2021}.
The value of insurance under the ZARC was positive at all suitable dates and highest for early sowing,
which reinforces the zoning as an incentive to sow early. The fixed premium used by the official cost survey
was higher than the actuarially fair premium estimated here, so the value of insurance to the farmer is
probably underestimated in the main scenario, even before the premium subsidy. The agreement of the modelled
loss frequency with the indemnified PSR policies between years ($r$ = 0.88) supports the use of the hybrid
yield distribution for risk analysis, although the extreme season of 2021 was underestimated.

The regret analysis adds a caveat on the economic value of more accurate yield models. Whether to sow
second-season maize in a given year depended mostly on price and cost: the realized margin of always sowing
was close to zero, and the improvement of 30--100 kg ha$^{-1}$ in $RMSE$ between models corresponded to
differences in regret within the uncertainty of the bootstrap. Better yield prediction thus does not
necessarily lead to better binary decisions when the outcome is dominated by economics. The yield model is
decisive, instead, where the physics separates the alternatives, as in the choice of the sowing date, in
which the date-specific water, frost and heat conditions determine the ranking. The ZARC replica illustrates
the same point from another angle: with a fixed rooting depth, the MZA-FAO water criterion agreed with the
official zoning in only 39\% of the combinations, so the official table, rather than the replica, was used
to define insurance access.

\subsection{Limitations}
\label{subsec:dis_limitations}

The analysis has limitations that should be considered when applying the results. Municipal yields do not
report sowing dates, so the hybrid model was trained on the mean of a sowing window and applied to individual
dates; the monotonic constraints and the consistency checks reduce, but do not eliminate, the uncertainty of
this extrapolation. Costs and prices were state averages, without regional differences in freight or price,
and costs did not vary with the sowing date. The opportunity cost of very early soybean cultivars and the
effects of pests and diseases were not represented. The thermal time of 1,450 \textdegree{}C d, taken from the
literature, may be short for current hybrids, since 1,550 \textdegree{}C d improved the hybrid model; it was
kept to avoid tuning the crop model on the validation data, and the best date was robust to this change.
Finally, the field-level risk was calibrated with a single parameter on nine years of insurance data, which
constrains its mean level but not its spatial pattern.

% =========================================================
% 5 CONCLUSION
% =========================================================
\section{Conclusion}
\label{sec:conclusion}

\begin{itemize}[leftmargin=*, itemsep=2pt]
\item The hybrid model of second-season maize yield explained 43\% of the variance of municipal yields,
54--61\% at the regional scales and 67\% of the variation of the state mean between years, validated on
seasons not used in training.
\item Uncalibrated APSIM, DSSAT and AquaCrop did not improve the hybrid model in relation to the simple
MZA-FAO; with an identical feature selection, MZA-FAO, APSIM and DSSAT were equivalent and AquaCrop was the
least informative.
\item Water satisfaction at flowering and grain filling, frost in the reproductive phase and heat at
flowering were the main drivers of second-season maize yield in Paraná.
\item Early sowing maximized the certainty equivalent of the margin in 84\% of the municipalities, with
the lowest probability of loss and the highest value of insurance under the ZARC; the warm lowlands of the
west and northwest preferred mid-March sowing because of heat at flowering.
\item The recommendation was robust to the cost basis, insurance terms, prices, crop model thermal time and
training design, and the modelled risk followed the indemnified insurance policies between years.
\end{itemize}

% =========================================================
% AUTHOR CONTRIBUTIONS
% =========================================================
\section*{Author Contributions}

F.C.O.: Conceptualization, Methodology, Software, Formal analysis, Data curation, Writing -- original draft,
Supervision. [INSERT: contributions of the co-authors (CRediT).]


% =========================================================
% ACKNOWLEDGEMENTS
% =========================================================
\section*{Acknowledgements}

The authors thank the Universidade Tecnológica Federal do Paraná (UTFPR), Campus Santa Helena, and
[INSERT: CAPES, Fundação Araucária or other supporting institutions].


% =========================================================
% FUNDING
% =========================================================
\section*{Funding}

[INSERT: funding sources and grant numbers, or ``This research received no specific grant from any
funding agency.'']


% =========================================================
% CONFLICT OF INTEREST
% =========================================================
\section*{Conflict of Interest}

The authors declare no conflicts of interest.


% =========================================================
% DATA AVAILABILITY
% =========================================================
\section*{Data Availability Statement}

The code, the processed data and the results that support the findings of this study are openly
available at \url{https://github.com/fcoliveira-utfpr/ML_maize_pr}. The daily weather data (BR-DWGD) are
available on Google Earth Engine, the municipal yields from the IBGE SIDRA system (table 839), the ZARC
risk tables and the crop insurance records from the open data portal of the Brazilian Ministry of
Agriculture, and the production costs and prices from SEAB/DERAL, CONAB and Ipeadata.


% =========================================================
% SUPPLEMENTARY INFORMATION
% =========================================================
\section*{Supplementary Information}

Supplementary Table S1, comparison of machine learning algorithms; Table S2, comparison of the feature
generators with an identical automatic feature selection; Table S3, regret of the plant/do-not-plant
decision for each yield model; Table S4, municipalities whose best sowing date was 15 March;
Figure S1, real production costs and prices of second-season maize; Figure S2, reliability of the
out-of-sample probability of loss (file \texttt{supplementary.pdf}).




% =========================================================
% REFERENCES
% =========================================================
\section*{References}
% Keep the reference list consistent. The journal uses APA style, but Free Format submission
% allows references to be submitted in another consistent style at initial submission.
\printbibliography[heading=none]


% =========================================================
% END DOCUMENT
% =========================================================
\end{document}